\subsection{Linear independence of homomorphisms}

Let L be an extension of K and V a vector space over K. In this paragraph we shall denote by \(\operatorname{Hom}_{\mathbf{K}}(\mathbf{V},\mathbf{L})\) the set of all K-linear mappings of V into L, equipped with the vector space structure on L such that :

\[
(f + g) (x) = f (x) + g (x), \quad (\alpha f) (x) = \alpha f (x)\tag{1}
\]

for \(x \in \mathbf{V}\) , \(\alpha \in \mathbf{L}\) and \(f, g\) in \(\operatorname{Hom}_{\mathbf{K}}(\mathbf{V}, \mathbf{L})\) . Let \(\mathbf{V}_{(\mathbf{L})} = \mathbf{L} \otimes_{\mathbf{K}} \mathbf{V}\) be the vector space on \(\mathbf{L}\) derived from \(\mathbf{V}\) by extension of scalars, and \((\mathbf{V}_{(\mathbf{L})})^*\) its dual. By II, p.~277 we have a canonical isomorphism \(u \mapsto \tilde{u}\) of vector \(\mathbf{L}\) -spaces from \((\mathbf{V}_{(\mathbf{L})})^*\) onto \(\operatorname{Hom}_{\mathbf{K}}(\mathbf{V}, \mathbf{L})\) such that \(\tilde{u}(x) = u(1 \otimes x)\) for \(x \in \mathbf{V}\) and \(u\) in \((\mathbf{V}_{(\mathbf{L})})^*\) . If \(\mathbf{V}\) is of finite dimension \(n\) over \(\mathbf{K}\) , the vector space \((\mathbf{V}_{(\mathbf{L})})\) over \(\mathbf{L}\) is of dimension \(n\) , as well as its dual \((\mathbf{V}_{(\mathbf{L})})^* = \mathbf{V}_{(\mathbf{L})}^*\) , whence the formula

\[
[ \operatorname{Hom} _ {K} (V, L): L ] = [ V: K ].\tag{2}
\]

THEOREM 1. --- Let L be an extension of a field K and A an algebra over K ; let H be the set of all K-algebra homomorphisms of A into L. Then H is a free subset of the vector space \(\operatorname{Hom}_{K}(A, L)\) over L.

Let us show by induction on the integer \(n \geqslant 0\) that every sequence \((u_1, \ldots, u_n)\) of distinct elements of \(\mathcal{H}\) is free. The case \(n = 0\) being trivial, we may henceforth suppose that \(n \geqslant 1\) ; let \(\alpha_1, \ldots, \alpha_n\) be elements of \(L\) such that \(\sum_{i=1}^{n} \alpha_i u_i = 0\) . For \(x, y\) in \(A\) we have

\[
\sum_ {i = 1} ^ {n - 1} \alpha_ {i} [ u _ {i} (x) - u _ {n} (x) ] \cdot u _ {i} (y) = \sum_ {i = 1} ^ {n} \alpha_ {i} u _ {i} (x y) - u _ {n} (x) \sum_ {i = 1} ^ {n} \alpha_ {i} u _ {i} (y) = 0,
\]

whence \(\sum_{i=1}^{n-1} \alpha_i [u_i(x) - u_n(x)] \cdot u_i = 0\) . By the induction hypothesis, the elements

\(u_{1}, \ldots, u_{n-1}\) of \(\mathscr{H}\) are linearly independent, whence \(\alpha_{i}[u_{i}(x) - u_{n}(x)] = 0\) for \(1 \leqslant i \leqslant n-1\) and for all x in A. Since the \(u_{i}\) are distinct, this implies that \(\alpha_{i} = 0\) for \(i \neq n\) , hence \(\alpha_{n}u_{n} = 0\) and so \(\alpha_{n} = \alpha_{n}u_{n}(1) = 0\) (on denoting by 1 the unit element of A). We have thus shown that \(\alpha_{1}, \ldots, \alpha_{n-1}, \alpha_{n}\) are zero, and this proves the theorem.

COROLLARY 1. --- Let \(\Gamma\) be a monoid, L a field and X a set of homomorphisms of \(\Gamma\) into the multiplicative monoid of L. Then X is a free subset of the vector L-space \(L^{\Gamma}\) of mappings of \(\Gamma\) into L.

Let A be the algebra of the monoid \(\Gamma\) with coefficients in L and \((e_{\gamma})_{\gamma \in \Gamma}\) the canonical basis of A over L (III, p.~446). For every L-linear mapping \(u\) of A into L let us write \(\tilde{u}(\gamma) = u(e_{\gamma})\) (for \(\gamma \in \Gamma\) ); then the mapping \(u \mapsto \tilde{u}\) is an isomorphism of vector L-spaces of \(\operatorname{Hom}_{L}(A, L)\) onto \(L^{\Gamma}\) which maps onto X the set of L-algebra homomorphisms of A into L. Now it suffices to apply Th. 1 with \(K = L\) .

COROLLARY 2 (Dedekind's theorem). --- Let E and L be two extensions of K. The set of K-homomorphisms of E into L is free over L. If E is of finite degree over K, the number of K-homomorphisms of E into L is at most equal to {[}E:K{]}.

The last assertion follows from the first, taking account of Formula (2).

